Modern cloud services, including web search, online retail, and serverless platforms, operate under microsecond-scale latency Service Level Objectives (SLOs) \cite{barroso2017attack,kaffes2019shinjuku,marty2019snap,zhang2021demikernel,iyer2023achieving}.
With every user request fanned out across thousands of servers, the end-to-end response time is determined by the slowest server.
Therefore, to ensure a satisfactory end-to-end responsiveness, bounding tail latency at each individual server is crucial.
These SLOs are becoming increasingly stringent, driven by the widespread deployment of ultra-low-latency hardware \yihan{cite} and the rise of microservice architectures \yihan{cite}.

Unfortunately, general-purpose operating systems (OSes) like Linux struggle to meet microsecond-scale tail latency SLOs due to unpredictable overheads in kernel I/O stack and scheduler.
To compensate, prior systems proposed \textit{kernel bypassing} to avoid OS intervention on the datapath, introducing the concept of the \textit{dataplane OS}.
Dataplane OSes move I/O processing to userspace and access network devices directly without involving the kernel.
Systems such as IX~\cite{belay2014ix} use hardware virtualization features like Intel VT-d~\cite{} to achieve protected, direct NIC queue access.
Other systems such as Shenango~\cite{} and Persephone~\cite{} rely on high-performance user-space I/O libraries such as the Data Plane Development Kit (DPDK)~\cite{}.
For example, IX~\cite{belay2014ix} uses hardware virtualization features such as Intel VT-d~\cite{} to safely expose NIC queues to userspace.
Other systems, like Shenango~\cite{} and Persephone~\cite{}, rely on user-space I/O libraries such as the Data Plane Development Kit (DPDK)~\cite{} to achieve similar goals.
Additionally, dataplane OSes implement a userspace scheduler to manage lightweight user-level threads.
IX adopts a single-threaded, run-to-completion model that reduces ring transition overhead through adaptive batching.
In contrast, Shenango and Arachne~\cite{} use cooperative multithreading, where threads yield control explicitly to the userspace scheduler.

Yet kernel bypassing alone is insufficient, as neither run-to-completion nor cooperative scheduling can prevent head-of-line (HoL) blocking under the high-dispersion workloads common in cloud services.
In such scenarios, a long-running task can monopolize a core, delaying the execution of short, latency-sensitive tasks and severely inflating tail latency.
To avoid violating tail latency SLOs, dataplane OSes must support \textit{preemptive scheduling}, allowing latency-critical tasks to receive timely execution.

Prior work on dataplane OS has explored various mechanisms for implementing preemptive scheduling in userspace.
% \yihan{For example, Shinjuku leverages unicast Interprocessor Interrupts (IPIs) to preempt tasks who have exhausted their 5$\mu$s timeslices.
% Concord further reduces preemption overhead by using cooperative preemption via cache-coherent message passing \cite{iyer2023achieving}.
% Signal?}

However, existing dataplane OSes typically rely on a \textit{centralized dispatcher} to preempt remote worker cores.
As core counts grow in modern servers, the cost of issuing these preemptions through a centralized dispatcher increases proportionally, creating a fundamental scalability bottleneck.
The system must then trade off preemption precision against resource efficiency:
it can either reduce preemption frequency, increasing the risk of tail latency SLO violations, or dedicate more cores to dispatching, reducing resources available to applications.
% As core counts in modern servers continue to increase, the centralized preemption dispatcher faces a fundamental scalability dilemma:
% the cost of issuing preemptions grows with the number of cores, eventually forcing the system to choose between sacrificing preemption granularity or dedicating more cores to maintain precision.
% This hidden scalability bottleneck limits the effectiveness of preemptive scheduling in large multi-core servers.
% it must balance preemption precision against preemption overhead.
% On one hand, issuing preemptions individually ensures microsecond-scale precision but incurs preemption overhead that grows linearly with core count. 
% On the other hand, batching preemptions reduces overhead but introduces timing skew across cores, leading to inconsistent timeslicing, and potential SLO violations.
% To mitigate this dilemma, the dataplane OS could reserve additional cores solely for preemption issuing to guarantee preemption precision.
% This creates a compounded scalability challenge: preserving preemption precision both increases preemption overhead and reduces core efficiency.

In this work, we address the scalability dilemma with \sysname, a new dataplane OS that enables distributed, microsecond-scale preemptive scheduling on multi-core servers.
Our key insight is that modern processors already provide hardware support for fine-grained, per-core preemption through the Local Advanced Programmable Interrupt Controller (LAPIC) timer.
However, the LAPIC timer is traditionally reserved for kernel use and is inaccessible by dataplane OS schedulers.
This leads to the central question of our work:
\textit{Can we safely expose kernel-critical hardware resources like the LAPIC timer to the dataplane OS scheduler?}

Prior dataplane OSes have explored isolation mechanisms in userspace, but they are not designed to securely multiplex hardware resources between the kernel and a userspace scheduler.
IX, ZygOS, and Arrakis expose I/O devices like NIC queues to applications using hardware virtualization (e.g., VT-x, IOMMU), assuming exclusive ownership and coarse-grained protection.
Vessel uses MPK to isolate a privileged userspace scheduler from untrusted application code within the same address space, but it does not expose or protect low-level hardware resources like LAPIC timer.

\sysname fills this gap by introducing privilege separation within userspace for managing kernel-critical hardware resources.
At the core of this design is the \textit{Shadow Ring}, a lightweight software enclave that shares the same privilege level (Ring 3) as application code but is memory-isolated using MPK.
Unlike prior systems that rely on full device ownership to userspace, \sysname enables secure multiplexing of kernel-critical hardware like the LAPIC timer to a trusted scheduler in userspace.
Shadow Ring allows \sysname to extend the Exokernel philosophy beyond I/O devices, enabling direct, low-latency control over system-critical resources without compromising safety.
The \sysname scheduler therefore operates across two isolated zones:
\begin{enumerate*}[label=(\roman*)]
  \item an untrusted zone that manages user-level threads and executes alongside application code, and
  \item a trusted zone (\schedcore), which runs in the Shadow Ring and interacts directly with privileged hardware.
\end{enumerate*}

Although the Shadow Ring isolates trusted scheduler logic in userspace, access to system-critical hardware still requires mediation from the kernel.
To support this, \sysname introduces a minimal kernel module (\kmod) that securely mediates access between \schedcore and hardware components such as the LAPIC timer.
To configure a hardware resource, \schedcore can either 
\begin{enumerate*}[label=(\arabic*)]
  \item issue a \textit{downcall} via a system call (e.g., \texttt{ioctl()}) that traps into Ring 0, where \kmod completes hardware configuration, or
  \item directly write to memory-mapped I/O (MMIO) regions of the hardware, which is mapped into a MPK-protected data segment in the Shadow Ring.
\end{enumerate*}
When a hardware event such as a timer expiration occurs, \kmod delivers it to the scheduler via an \textit{upcall}, transferring control to a registered handler located in the MPK-protected text segment in the Shadow Ring.

We compare \sysname with Shinjuku~\cite{kaffes2019shinjuku}, Perséphone~\cite{demoulin2021idling} and Concord~\cite{iyer2023achieving}, three state-of-the-art dataplane operating systems.
For synthetic workloads, we show that \yihan{xx}.
For YCSB workloads on RocksDB, a widely used key-value store, we show that \yihan{xx}.
We show that \sysname scales better than prior systems as core counts increase, while preserving microsecond-scale preemption precision.

The rest of this paper is organized as follows.
\S~\ref{sec:background} motivates our work by discussing the scalability challenges of existing preemptive scheduling systems.
\S~\ref{sec:challenge} discusses the key challenges and design approaches of \sysname.
\S~\ref{sec:design} presents the detailed design and implementation of \sysname.